\section{Gestión de Riesgos}
\label{sec:gestion_riesgos}

%La gestión de riesgos considera las condiciones que pueden afectar la interacción entre el Sistema Coordinador, los servidores MySQL y las distintas instancias de OCR Server.

\subsection{Supuestos}
\label{subsec:supuestos}

\begin{enumerate}
    \item Los servidores MySQL se encuentran disponibles y permiten el acceso de usuarios autorizados.
    \item Las tablas seleccionadas poseen una clave primaria y una columna que permite recuperar el documento.
    \item Los OCR Server informan su disponibilidad, modelo activo y capacidad de procesamiento.
    \item Existe conectividad entre el Sistema Coordinador, MySQL y los OCR Server.
    \item Los documentos poseen formatos compatibles con alguno de los perfiles disponibles.
    \item SQLite permite conservar el catálogo estructural y el estado de las tareas.
\end{enumerate}

\subsection{Dependencias}
\label{subsec:dependencias}

\begin{enumerate}
    \item \textbf{MySQL institucional:} proporciona los documentos, metadatos y estructuras de destino.
    \item \textbf{OCR Server:} ejecuta el reconocimiento y análisis documental.
    \item \textbf{Red institucional:} permite la comunicación entre los componentes.
    \item \textbf{Recursos computacionales:} determinan la capacidad y velocidad de procesamiento.
    \item \textbf{Modelos y bibliotecas:} permiten ejecutar los perfiles OCR configurados.
    \item \textbf{SQLite:} mantiene el catálogo local, los servidores registrados y las tareas.
\end{enumerate}

\subsection{Restricciones}
\label{subsec:restricciones}

\begin{enumerate}
    \item El sistema deberá funcionar con el hardware disponible en la institución.
    \item Cada OCR Server deberá respetar su capacidad y nivel de concurrencia.
    \item Una tarea solo podrá asignarse a servidores compatibles con el modelo requerido.
    \item Los documentos y registros originales no deberán eliminarse ni sobrescribirse.
    \item Toda ejecución, reasignación y reprocesamiento deberá conservar su trazabilidad.
    \item SQLite almacenará información estructural y de control, pero no reemplazará a MySQL.
    \item Los componentes internos de OCR Server deberán permanecer dentro de una red controlada.
\end{enumerate}

\newpage

\subsection{Riesgos}
\label{subsec:riesgos}

La Tabla~\ref{tab:riesgos_sistema} presenta los principales riesgos identificados y sus medidas de mitigación.

\renewcommand{\arraystretch}{1.2}

\begin{longtable}{
    >{\raggedright\arraybackslash}p{0.43\textwidth}
    >{\raggedright\arraybackslash}p{0.49\textwidth}
}
\caption{Riesgos del sistema y medidas de mitigación.}
\label{tab:riesgos_sistema}\\

\toprule
\textbf{Riesgo} & \textbf{Medida de mitigación} \\
\midrule
\endfirsthead

\toprule
\textbf{Riesgo} & \textbf{Medida de mitigación} \\
\midrule
\endhead

\bottomrule
\endfoot

\textbf{Indisponibilidad de un OCR Server.}

Una instancia puede dejar de responder durante una tarea.
&
Detectar su estado y reasignar los documentos pendientes a otro servidor compatible.
\\

\textbf{Saturación de recursos.}

La carga puede superar la memoria o capacidad de procesamiento disponible.
&
Limitar la concurrencia y distribuir las tareas según la carga de cada servidor.
\\

\textbf{Asignación incompatible.}

Una tarea puede enviarse a un servidor que no dispone del modelo requerido.
&
Validar previamente los modelos, recursos y capacidades de cada instancia.
\\

\textbf{Pérdida o duplicación de tareas.}

Una interrupción puede provocar ejecuciones incompletas o repetidas.
&
Utilizar identificadores únicos, estados persistentes y límites de reintentos.
\\

\textbf{Desactualización del catálogo.}

La estructura de MySQL puede cambiar respecto de la registrada en SQLite.
&
Sincronizar y validar el catálogo antes de crear o ejecutar una tarea.
\\

\textbf{Resultados OCR incorrectos.}

La calidad del documento o del modelo puede afectar la transcripción.
&
Conservar el original y permitir el reprocesamiento con otro modelo o perfil.
\\

\textbf{Fallo al almacenar resultados.}

Una escritura incompleta puede generar inconsistencias.
&
Utilizar transacciones y confirmar la persistencia antes de finalizar el documento.
\\

\textbf{Exposición de información.}

Credenciales o servicios internos pueden ser accedidos sin autorización.
&
Restringir permisos, proteger las sesiones y operar dentro de una red controlada.
\\

\textbf{Crecimiento del almacenamiento.}

Los archivos temporales, resultados y registros pueden consumir el espacio disponible.
&
Aplicar políticas de limpieza, respaldo y retención de información.
\\

\end{longtable}